% Source: 27.quadratic forms.v03, twelve pages numbering 27.a.1 through 27.a.9.
% The Polarization Identity carries no number in the notes (Prof. Biran marks
% it with a circled star), hence theorem* here.
\renewcommand{\thetheorem}{27.a.\arabic{theorem}}
\setcounter{theorem}{0}

\setcounter{chapter}{26}
% ---------------------------------------------------------------------------
% Provenance of the prose in this file.
%   % Extractor: <model>   the block below was transcribed from Prof. Biran's
%                          handwritten notes by that model
%   % Creator:   <model>   the block below was written by that model and has no
%                          counterpart in the notes
% A tag holds until the next tag.
% ---------------------------------------------------------------------------

% Extractor: Gemini 3.1 Pro
\chapter{Bilinear and Quadratic Forms}

% Def. (27.a.1)
\begin{definition}[Bilinear Form]
\label{def:bilinear_form}
Let $V$ be a vector space over $K$ (a general field). A function $B \colon V \times V \to K$ is called a \newterm{bilinear form} if the following two conditions are satisfied:
\begin{enumerate}[label=\textbf{(\alph*)}]
    \item For all $w \in V$, the map $B(-, w) \colon V \to K$, $v \mapsto B(v, w)$ is linear.
    \item For all $v \in V$, the map $B(v, -) \colon V \to K$, $w \mapsto B(v, w)$ is linear.
\end{enumerate}
\end{definition}

% Lemma. (27.a.2)
\begin{lemma}[Matrix Representation of a Bilinear Form]
\label{lem:matrix_rep_bilinear_form}
Let $V$ be a vector space over $K$ and $\mathcal{C} := (e_1, \dots, e_n)$ a basis for $V$. Let $B \colon V \times V \to K$ be a bilinear form. Then there exists a matrix $[B]_{\mathcal{C}} \in \M_{n \times n}(K)$ such that
\[
B\left(\sum_{i=1}^n c_i e_i, \sum_{j=1}^n d_j e_j\right) = (c_1, \dots, c_n) \cdot [B]_{\mathcal{C}} \cdot (d_1, \dots, d_n)\transp \quad \forall c_i, d_j \in K.
\]
The matrix $[B]_{\mathcal{C}}$ is called the \newterm{matrix representation} of $B$ in the basis $\mathcal{C}$. (It is unique, which is what entitles us to the definite article; see the end of the proof).
\end{lemma}

\begin{proof}
Define $[B]_{\mathcal{C}} := (a_{ij})$, where $a_{ij} := B(e_i, e_j)$. Then we compute:
\begin{align*}
B\left(\sum_{i=1}^n c_i e_i, \sum_{j=1}^n d_j e_j\right) &= \sum_{i=1}^n c_i B\left(e_i, \sum_{j=1}^n d_j e_j\right) \\
&= \sum_{i=1}^n c_i \left(\sum_{j=1}^n d_j B(e_i, e_j)\right) \\
&= \sum_{i=1}^n \sum_{j=1}^n c_i d_j B(e_i, e_j) \\
&= (c_1, \dots, c_n) \cdot [B]_{\mathcal{C}} \cdot \begin{pmatrix} d_1 \\ \vdots \\ d_n \end{pmatrix}.
\end{align*}
% Creator: Opus 5
For uniqueness, suppose $M \in \M_{n \times n}(K)$ also satisfies the displayed identity. Fix $i$ and $j$ and choose the coefficients $c_1, \dots, c_n$ to be $1$ at the index $i$ and $0$ elsewhere, and $d_1, \dots, d_n$ to be $1$ at the index $j$ and $0$ elsewhere. The left-hand side is then $B(e_i, e_j)$ and the right-hand side is the single entry $M_{ij}$, so $M_{ij} = B(e_i, e_j) = ([B]_{\mathcal{C}})_{ij}$ for all $i$ and $j$, that is, $M = [B]_{\mathcal{C}}$.
% Extractor: Gemini 3.1 Pro
\end{proof}

We saw that also for $T \in \End(V)$ and $\mathcal{C}$ a basis for $V$, we have a matrix representation $[T]_{\mathcal{C}}^{\mathcal{C}} \in \M_{n \times n}(K)$, and that if $\mathcal{D}$ is another basis for $V$, then
\begin{equation}
\label{eq:transition_endomorphism}
[T]_{\mathcal{D}}^{\mathcal{D}} = [\id_V]_{\mathcal{D}}^{\mathcal{C}} \cdot [T]_{\mathcal{C}}^{\mathcal{C}} \cdot [\id_V]_{\mathcal{C}}^{\mathcal{D}},
\end{equation}
where $[\id_V]_{\mathcal{C}}^{\mathcal{D}} = ([\id_V]_{\mathcal{D}}^{\mathcal{C}})^{-1}$ is the transition matrix from basis $\mathcal{D}$ to $\mathcal{C}$.

\begin{question}
What happens for bilinear forms?
\end{question}

% Lemma. (27.a.3)
\begin{lemma}[Change of Basis for Bilinear Forms]
\label{lem:change_of_basis_bilinear}
Let $V$ be a finite-dimensional vector space over $K$ and $\mathcal{C}, \mathcal{D}$ two bases of $V$. Let $B \colon V \times V \to K$ be a bilinear form. Then
\[
[B]_{\mathcal{D}} = \left([\id_V]_{\mathcal{C}}^{\mathcal{D}}\right)\transp \cdot [B]_{\mathcal{C}} \cdot [\id_V]_{\mathcal{C}}^{\mathcal{D}}.
\]
(Note that this is quite different from \eqref{eq:transition_endomorphism}).
\end{lemma}

\begin{proof}
Let $\mathcal{C} := (v_1, \dots, v_n)$ and $\mathcal{D} := (w_1, \dots, w_n)$. We have $([B]_{\mathcal{C}})_{j,k} = B(v_j, v_k)$, and $([B]_{\mathcal{D}})_{j,k} = B(w_j, w_k)$ (this is how the previous proof constructed those matrices, and by the uniqueness just established there is no other reading). Recall that column $j$ in $[\id_V]_{\mathcal{C}}^{\mathcal{D}}$ contains the coordinates of $w_j$ in the basis $\mathcal{C}$, by the Proposition on \cpageref{prop:columns_of_representation_matrix} applied to $\id_V$:
\[
w_j = \sum_{p=1}^n \left([\id_V]_{\mathcal{C}}^{\mathcal{D}}\right)_{p,j} \cdot v_p.
\]
This implies:
\begin{align*}
([B]_{\mathcal{D}})_{j,k} &= B(w_j, w_k) \\
&= B\left(\sum_{p=1}^n \left([\id_V]_{\mathcal{C}}^{\mathcal{D}}\right)_{p,j} v_p, \sum_{q=1}^n \left([\id_V]_{\mathcal{C}}^{\mathcal{D}}\right)_{q,k} v_q\right) \\
&= \sum_{p,q=1}^n \left([\id_V]_{\mathcal{C}}^{\mathcal{D}}\right)_{p,j} \cdot B(v_p, v_q) \cdot \left([\id_V]_{\mathcal{C}}^{\mathcal{D}}\right)_{q,k} \\
&= \left(\left([\id_V]_{\mathcal{C}}^{\mathcal{D}}\right)\transp \cdot [B]_{\mathcal{C}} \cdot [\id_V]_{\mathcal{C}}^{\mathcal{D}}\right)_{j,k}.
\end{align*}
\end{proof}

% Def. (27.a.4)
\begin{definition}[Symmetric, Positive Definite, and Quadratic Forms]
\label{def:symmetric_positive_definite_quadratic}
\leavevmode
\begin{enumerate}[label=\textbf{(\alph*)}]
    \item A bilinear form $B \colon V \times V \to K$ is called \newterm{symmetric} if $B(v, w) = B(w, v)$ for all $v, w \in V$.
    \item Assume $K = \mathbb{R}$. A bilinear form $B \colon V \times V \to \mathbb{R}$ is called \newterm{positive definite} if $B(v, v) > 0$ for all $0 \neq v \in V$.
    \item For a bilinear form $B \colon V \times V \to K$, we define $q_B \colon V \to K$ by $q_B(v) := B(v, v)$. We call $q_B$ the \newterm{quadratic form} associated to $B$. Sometimes we just denote it by $q$.
\end{enumerate}
\end{definition}

\begin{example}[Standard Bilinear and Quadratic Forms]
\leavevmode
\begin{enumerate}
    \item Let $(V, \langle \cdot, \cdot \rangle)$ be a Euclidean vector space. Then $B(v, w) := \langle v, w \rangle$ is a symmetric, positive definite bilinear form. Its associated quadratic form is $q(v) = \|v\|^2$.
    \item Let $B \in \M_{n \times n}(K)$. Then $B \colon K^n \times K^n \to K$ defined by $B(v, w) := v\transp B w$ is a bilinear form. (The letter $B$ is doing double duty here, standing both for the matrix and for the form it defines; throughout this item, $B$ written with two arguments is the form and $B$ written alone is the matrix). If $B\transp = B$ (i.e., $B$ is a symmetric matrix), then $B$ is a symmetric bilinear form (left as an exercise, see \cref{exc:symmetric_matrix_symmetric_form}). The associated quadratic form is $q(v) = v\transp B v$.
\end{enumerate}
\end{example}

% Extractor: Opus 5
\begin{exercise}[A Symmetric Matrix Gives a Symmetric Bilinear Form]
\label{exc:symmetric_matrix_symmetric_form}
This is the step Prof. Biran marks with \qt{exc.} in the example above. Let $B \in \M_{n \times n}(K)$ with $B\transp = B$, and let $B \colon K^n \times K^n \to K$ be the bilinear form $B(v, w) := v\transp B w$. Show that it is symmetric, that is, $B(v, w) = B(w, v)$ for all $v, w \in K^n$.
\exinfo{This exercise is taken from Prof.~Biran's handwritten lecture notes.}
\exsol{sol:symmetric_matrix_symmetric_form}
\end{exercise}
% Extractor: Gemini 3.1 Pro

\begin{remark}
Note that $B \in \M_{n \times n}(K)$ and $B\transp \in \M_{n \times n}(K)$ give the same quadratic form $q \colon K^n \to K$. Indeed, for all $v \in K^n$, we have:
\[
q_B(v) = v\transp B v = (v\transp B v)\transp = v\transp B\transp v = q_{B\transp}(v),
\]
where we used the fact that $v\transp B v$ is a scalar, and thus equals its own transpose.

Assume now that $2 \neq 0$ in $K$ (i.e., $1 + 1 \neq 0$). Then $\frac{1}{2}(B + B\transp)$ also gives the same quadratic form as $B$.
% Creator: Opus 5
Indeed, for all $v \in K^n$,
\[
q_{\frac{1}{2}(B + B\transp)}(v) = \tfrac{1}{2}\left( v\transp B v + v\transp B\transp v \right) = \tfrac{1}{2}\left( q_B(v) + q_{B\transp}(v) \right) = q_B(v),
\]
the last step by the identity just proved. What makes this worth recording is not the equality itself but the symmetry of the new matrix: $\left(\frac{1}{2}(B + B\transp)\right)\transp = \frac{1}{2}(B\transp + B)$, so every quadratic form on $K^n$ is $q_A$ for a \emph{symmetric} $A$, which is exactly what the Conclusion below asserts.
% Extractor: Gemini 3.1 Pro
\end{remark}

\begin{conclusion}
If $2 \neq 0$ in $K$, then for the discussion on quadratic forms $q \colon K^n \to K$, we can assume they all come from symmetric matrices (or symmetric bilinear forms).
\end{conclusion}

\begin{theorem*}[Polarization Identity]
\label{thm:polarization_identity_bilinear}
Assume $2 \neq 0$ in $K$. (Note that $2 \neq 0 \implies 4 \neq 0$. Why? This is left as an exercise, see \cref{exc:two_nonzero_four_nonzero}).
If $B$ is symmetric and defines $q := q_B$, then for all $v, w \in V$ we have:
\begin{enumerate}[label=\textbf{(\alph*)}]
    \item $B(v, w) = \frac{1}{4}q(v + w) - \frac{1}{4}q(v - w)$.
    \item $B(v, w) = \frac{1}{2}(q(v + w) - q(v) - q(w))$.
\end{enumerate}
\end{theorem*}

% Creator: Opus 5
% The notes state the identity with no proof; it is two expansions, and the
% chapter uses it at a decisive point (the base case of 27.a.7), so it is worth
% having on the page.
\begin{proof}
Bilinearity and the symmetry of $B$ give, for all $v, w \in V$,
\[
    q(v \pm w) = B(v \pm w, v \pm w) = B(v, v) \pm B(v, w) \pm B(w, v) + B(w, w) = q(v) \pm 2 B(v, w) + q(w).
\]
Subtracting the two versions of this identity gives $q(v+w) - q(v-w) = 4 B(v, w)$, which is \textbf{(a)}, division by $4$ being legitimate by \cref{exc:two_nonzero_four_nonzero}. Taking only the version with the plus sign and moving $q(v)$ and $q(w)$ to the left gives $q(v+w) - q(v) - q(w) = 2 B(v, w)$, which is \textbf{(b)}.
\end{proof}
% Extractor: Gemini 3.1 Pro

% Extractor: Opus 5
\begin{exercise}[Invertibility of 2 and 4 in Characteristic \texorpdfstring{$\neq 2$}{\ne 2}]
\label{exc:two_nonzero_four_nonzero}
This is the question Prof. Biran annotates the Polarization Identity with. Let $K$ be a field in which $2 \neq 0$. Show that then also $4 \neq 0$, so that the divisions by $4$ in part \textbf{(a)} really are legitimate.
\exinfo{This exercise is taken from Prof.~Biran's handwritten lecture notes.}
\exsol{sol:two_nonzero_four_nonzero}
\end{exercise}

\begin{exercise}[Polarization Identity for Quadratic Forms]
  \label{exc:polarization_quadratic_forms_series26}
  Let $K$ be a field with $\operatorname{char}(K) \neq 2$, let $V$ be a vector space over $K$, and let $B \colon V \times V \to K$ be a symmetric bilinear form on $V$. Let $q_B(v) := B(v, v)$ be the associated quadratic form. Prove that for all $v, w \in V$,
  \[
  B(v, w) = \frac{1}{2} \big( q_B(v + w) - q_B(v) - q_B(w) \big).
  \]
  \exhint{Expand $q_B(v + w) = B(v + w, v + w) = B(v, v) + B(v, w) + B(w, v) + B(w, w)$ using bilinearity. By symmetry, $B(w, v) = B(v, w)$, so $q_B(v + w) = q_B(v) + 2B(v, w) + q_B(w)$. Subtracting $q_B(v) + q_B(w)$ from both sides and dividing by $2$ (which is valid since $\operatorname{char}(K) \neq 2$) gives the result.}
  \exinfo{This exercise is Problem 1 of Problem Sheet 26, Linear Algebra II, Spring Semester 2026, ETH Zurich (Lecturer: Prof. P. Biran).}
\exsol{sol:polarization_quadratic_forms_series26}
\end{exercise}
% Extractor: Gemini 3.1 Pro

% Thm. (27.a.5)
\begin{theorem}[Sylvester's Law of Inertia]
\label{thm:sylvesters_inertia}
Let $B \colon V \times V \to \mathbb{R}$ be a \emph{symmetric} bilinear form on a vector space $V$ over $\mathbb{R}$, of dimension $n$. Then there exists a basis $\mathcal{C}$ for $V$ such that
\[
[B]_{\mathcal{C}} = \begin{pmatrix} I_k & 0 & 0 \\ 0 & -I_\ell & 0 \\ 0 & 0 & 0 \end{pmatrix},
\]
for some $k, \ell \in \mathbb{Z}_{\ge 0}$ with $k + \ell \le n$. (The zero block in the bottom right has size $n - (k + \ell)$).

Moreover, $k$ and $\ell$ do \emph{not} depend on the basis $\mathcal{C}$ (they depend only on $B$). In fact:
\begin{align*}
k &= \max \{ \dim W \mid W \subseteq V \text{ is a subspace and } B|_{W \times W} \text{ is positive definite} \}, \\
\ell &= \max \{ \dim U \mid U \subseteq V \text{ is a subspace and } -B|_{U \times U} \text{ is positive definite} \}, \\
n - (k + \ell) &= \max \{ \dim Z \mid Z \subseteq V \text{ is a subspace and } B|_{Z \times Z} \equiv 0 \}.
\end{align*}
\end{theorem}

\begin{remark}
\leavevmode
\begin{enumerate}[label=\textbf{(\alph*)}]
    \item In the basis $\mathcal{C}$, the quadratic form $q$ associated to $B$ looks like 
    \[ q(x_1, \dots, x_n) = x_1^2 + \dots + x_k^2 - x_{k+1}^2 - \dots - x_{k+\ell}^2. \]
    \item $k$ is called the \newterm{positive index} of $B$, $\ell$ is the \newterm{negative index}, and $n - (k + \ell)$ is the \newterm{nullity}. The value $\sigma(B) := \sgn(B) := k - \ell$ is called the \newterm{signature} of $B$. (This is sometimes referred to as type $(k, \ell, 0)$).
\end{enumerate}
\end{remark}

\begin{proof}
Choose any basis $\mathcal{C}'$ for $V$ and let $A := [B]_{\mathcal{C}'} \in \M_{n \times n}(\mathbb{R})$. Note that $A$ is a symmetric matrix (since $B$ is symmetric: $A_{ij} = B(v'_i, v'_j) = B(v'_j, v'_i) = A_{ji}$). By the Spectral Theorem for matrices (\cref{thm:spectral_thm_matrices} \textbf{(b)}), this implies that $A$ is orthogonally diagonalizable, i.e., there exists $P \in \Orth(n)$ such that $P^{-1} A P$ is a diagonal matrix. But $P^{-1} = P\transp$ (by \cref{lem:char_orthogonal_matrices} \textbf{(d)}), so $P\transp A P$ is diagonal.

Let $\mathcal{C}''$ be the (unique) basis of $V$ such that $[\id_V]_{\mathcal{C}'}^{\mathcal{C}''} = P$. (That is, if $\mathcal{C}' = (v'_1, \dots, v'_n)$ and $\mathcal{C}'' = (v''_1, \dots, v''_n)$, then $v''_j = \sum_{i=1}^n P_{ij} v'_i$ for all $j$).
We have, by \cref{lem:change_of_basis_bilinear} applied to the pair $\mathcal{C}'$, $\mathcal{C}''$,
\[
[B]_{\mathcal{C}''} = P\transp A P = \begin{pmatrix} \eta_1 & & 0 \\ & \ddots & \\ 0 & & \eta_n \end{pmatrix},
\]
where $\eta_j \in \mathbb{R}$ for all $j$.

By changing the order of the elements in the basis $\mathcal{C}''$, we may assume $\eta_1, \dots, \eta_k > 0$, $\eta_{k+1}, \dots, \eta_{k+\ell} < 0$, and $\eta_{k+\ell+1} = \dots = \eta_n = 0$ for some $k, \ell \in \mathbb{Z}_{\ge 0}$ with $k + \ell \le n$.

We can write:
\begin{align*}
\eta_i &= \lambda_i^2, \text{ with } \lambda_i > 0, \quad \forall 1 \le i \le k \quad (\text{so } B(v''_i, v''_i) = \lambda_i^2), \\
\eta_i &= -\lambda_i^2, \text{ with } \lambda_i > 0, \quad \forall k+1 \le i \le k+\ell \quad (\text{so } B(v''_i, v''_i) = -\lambda_i^2), \\
\eta_i &= 0, \quad \forall k+\ell+1 \le i \le n \quad (\text{so } B(v''_i, v''_i) = 0).
\end{align*}

Define $\mathcal{C} := (v_1, \dots, v_k, v_{k+1}, \dots, v_{k+\ell}, v_{k+\ell+1}, \dots, v_n)$, where 
\[
v_i = \frac{1}{\lambda_i} v''_i \quad \forall 1 \le i \le k+\ell, \qquad v_i = v''_i \quad \forall k+\ell+1 \le i \le n.
\]
In other words: 
\[
\mathcal{C} = \left(\frac{1}{\lambda_1} v''_1, \dots, \frac{1}{\lambda_k} v''_k, \frac{1}{\lambda_{k+1}} v''_{k+1}, \dots, \frac{1}{\lambda_{k+\ell}} v''_{k+\ell}, v''_{k+\ell+1}, \dots, v''_n\right).
\]
Note that $v_i = c_i v''_i$ for some scalars $c_i \neq 0$, for all $i$. It is left as an exercise to show that $\mathcal{C}$ is indeed a basis, see \cref{exc:scaled_family_is_basis}.

We have $([B]_{\mathcal{C}})_{i,j} = B(v_i, v_j) = B(c_i v''_i, c_j v''_j) = 0$ for all $i \neq j$.
\noindent
And for $i = j$:
\begin{itemize}
    \item if $1 \le i \le k$, then $([B]_{\mathcal{C}})_{i,i} = B(v_i, v_i) = B\left(\frac{1}{\lambda_i} v''_i, \frac{1}{\lambda_i} v''_i\right) = \frac{1}{\lambda_i^2} \lambda_i^2 = 1$,
    \item if $k+1 \le i \le k+\ell$, then $([B]_{\mathcal{C}})_{i,i} = B(v_i, v_i) = B\left(\frac{1}{\lambda_i} v''_i, \frac{1}{\lambda_i} v''_i\right) = \frac{1}{\lambda_i^2} (-\lambda_i^2) = -1$,
    \item if $k+\ell+1 \le i \le n$, then $([B]_{\mathcal{C}})_{i,i} = B(v_i, v_i) = B(v''_i, v''_i) = 0$.
\end{itemize}
This proves the existence of the basis $\mathcal{C}$ claimed in the theorem.

We will prove now that $k$ is independent of $\mathcal{C}$. Put 
\[
k' := \max \{ \dim W \mid W \subseteq V \text{ is a subspace such that } B|_{W \times W} \text{ is positive definite} \}.
\]

We will show $k = k'$. Indeed, $k \le k'$, because $B|_{U \times U}$ is positive definite on $U := \Sp(v_1, \dots, v_k)$. This is because:
\begin{align*}
    & B(a_1 v_1 + \dots + a_k v_k, a_1 v_1 + \dots + a_k v_k) \\
    & \quad = (a_1, \dots, a_k, 0, \dots, 0) \cdot \begin{pmatrix} I_k & 0 & 0 \\ 0 & -I_\ell & 0 \\ 0 & 0 & 0 \end{pmatrix} \begin{pmatrix} a_1 \\ \vdots \\ a_k \\ 0 \\ \vdots \\ 0 \end{pmatrix} \\
    & \quad = a_1^2 + \dots + a_k^2 \\
    & \quad > 0,
\end{align*}
if not all $a_i = 0$.

Suppose by contradiction that $k' \ge k+1$. This implies there exists $W \subseteq V$ of dimension $\ge k+1$ such that $B|_{W \times W}$ is positive definite. Take $X := \Sp(v_{k+1}, \dots, v_n)$.
Then, by the dimension formula (\cref{prop:properties_of_sums} \textbf{(c)}):
\[
\dim(W \cap X) = \dim W + \dim X - \dim(W+X) \ge (k+1) + (n-k) - \dim(W+X) \ge n+1 - n = 1.
\]
This implies $W \cap X \neq \{0\}$.
Take $0 \neq w \in W \cap X$. We can write $w = a_{k+1} v_{k+1} + \dots + a_n v_n$.
Then $B(w, w) = -a_{k+1}^2 - \dots - a_{k+\ell}^2 \le 0$. But, by the assumption that $W$ is positive definite, $B(w, w) > 0$. This is a contradiction, which implies $k' \le k$.
This proves $k = k'$.

The proofs for $\ell$ and $n - (k + \ell)$ are similar.
\end{proof}

% Extractor: Opus 5
\begin{exercise}[Rescaling a Basis Leaves a Basis]
\label{exc:scaled_family_is_basis}
This is the step marked \qt{Exc.} in the proof above. Let $\mathcal{C}'' := (v''_1, \dots, v''_n)$ be a basis for $V$, let $c_1, \dots, c_n \in K \setminus \{0\}$, and put $v_i := c_i v''_i$. Show that $\mathcal{C} := (v_1, \dots, v_n)$ is again a basis for $V$.
\exinfo{This exercise is taken from Prof.~Biran's handwritten lecture notes.}
\exsol{sol:scaled_family_is_basis}
\end{exercise}

\begin{exercise}[Why Symmetry Cannot Be Dropped in \texorpdfstring{\cref{thm:sylvesters_inertia}}{Sylvester's Law of Inertia}]
\label{exc:sylvester_needs_symmetry}
Prof. Biran annotates the word \qt{symmetric} in the statement of \cref{thm:sylvesters_inertia} with \qt{explain}. Explain why the hypothesis cannot simply be dropped: show that if $[B]_{\mathcal{C}}$ has the form displayed in the theorem for \emph{some} basis $\mathcal{C}$, then $B$ is necessarily symmetric. (Hint: that displayed matrix is symmetric, and \cref{lem:change_of_basis_bilinear} tells you what happens to it under a change of basis). Then give a bilinear form on $\mathbb{R}^2$ for which the conclusion of the theorem fails.
\exinfo{This exercise is taken from Prof.~Biran's handwritten lecture notes.}
\exsol{sol:sylvester_needs_symmetry}
\end{exercise}
% Extractor: Gemini 3.1 Pro

% Thm. (27.a.6)
\begin{theorem}[Sylvester's Law for Euclidean Spaces]
\label{thm:sylvesters_euclidean}
Let $V$ be a finite-dimensional Euclidean space of dimension $n$. Let $B \colon V \times V \to \mathbb{R}$ be a symmetric bilinear form. Then there exists an \emph{orthonormal} basis $\mathcal{C}$ for $V$ in which
\[
[B]_{\mathcal{C}} = \begin{pmatrix}
\eta_1 & & & & & & & & \\
& \ddots & & & & & & & \\
& & \eta_k & & & & & & \\
& & & \eta_{k+1} & & & & & \\
& & & & \ddots & & & & \\
& & & & & \eta_{k+\ell} & & & \\
& & & & & & 0 & & \\
& & & & & & & \ddots & \\
& & & & & & & & 0
\end{pmatrix},
\]
where $k + \ell \le n$, with $\eta_1, \dots, \eta_k > 0$ and $\eta_{k+1}, \dots, \eta_{k+\ell} < 0$, and where the remaining $n - (k+\ell)$ diagonal entries vanish. As before, $k, \ell$ are the positive and negative indices of $B$. Moreover, the scalars $\eta_1, \dots, \eta_k, \eta_{k+1}, \dots, \eta_{k+\ell}, 0, \dots, 0$ are independent of $\mathcal{C}$. They depend only on $B$ (and on the scalar product on $V$). They are the eigenvalues of $[B]_{\mathcal{D}}$ for every \emph{orthonormal} basis $\mathcal{D}$ of $V$ (independent of $\mathcal{D}$ and on whether or not $[B]_{\mathcal{D}}$ is diagonal).
\end{theorem}

\begin{remark}
This theorem has a geometric meaning.
If $q \colon \mathbb{R}^n \to \mathbb{R}$ is a quadratic form $q(x_1, \dots, x_n) = \sum a_{ij} x_i x_j$ for some $a_{ij} \in \mathbb{R}$, then the set
\[
Q = \{ (x_1, \dots, x_n) \in \mathbb{R}^n \mid q(x_1, \dots, x_n) = C \} \quad (\text{for some constant } C \in \mathbb{R})
\]
is called a \newterm{quadric}. (Examples include an ellipse/ellipsoid, parabola/paraboloid, hyperbola/hyperboloid).
The theorem tells us that there exists a linear isometry $T \colon \mathbb{R}^n \to \mathbb{R}^n$ such that
\[
T(Q) = \left\{ (y_1, \dots, y_n) \in \mathbb{R}^n \;\middle|\; \underbrace{\eta_1 y_1^2 + \dots + \eta_k y_k^2}_{\text{positive coeffs.}} + \underbrace{\eta_{k+1} y_{k+1}^2 + \dots + \eta_{k+\ell} y_{k+\ell}^2}_{\text{negative coeffs.}} = C \right\}.
\]
\end{remark}

\begin{proof}[Proof Outline]
Pick any orthonormal basis $\mathcal{C}'$ for $V$ (similar to the previous proof, only now we take $\mathcal{C}'$ to be orthonormal). Define $A := [B]_{\mathcal{C}'}$, and proceed exactly as in the previous proof, but instead of defining $\mathcal{C}$ with the scaling factors $\lambda_i$, simply take $\mathcal{C} = \mathcal{C}''$.

To see why $\eta_1, \dots, \eta_{k+\ell}, 0, \dots, 0$ depend only on $B$ etc., note that if $\mathcal{D}$ and $\mathcal{D}'$ are two orthonormal bases for $V$, then $[\id_V]_{\mathcal{D}'}^{\mathcal{D}} \in \Orth(n)$ (left as an exercise, see \cref{exc:transition_between_orthonormal_bases}). Put $Q := [\id_V]_{\mathcal{D}'}^{\mathcal{D}}$.
We have
\[
[B]_{\mathcal{D}} = Q\transp \cdot [B]_{\mathcal{D}'} \cdot Q = Q^{-1} [B]_{\mathcal{D}'} Q,
\]
since $Q \in \Orth(n)$, so that $Q\transp = Q^{-1}$ by \cref{lem:char_orthogonal_matrices} \textbf{(d)}. Hence, the matrices $[B]_{\mathcal{D}}$ and $[B]_{\mathcal{D}'}$ are similar, which means their eigenvalues are exactly the same: similar matrices share a characteristic polynomial by \cref{lem:properties_char_poly} \textbf{(e)}, and the eigenvalues are its roots by \cref{aiprop:eigenvalue_iff_root}. This is precisely where the congruence $Q\transp [B]_{\mathcal{D}'} Q$ of \cref{lem:change_of_basis_bilinear}, which in general does \emph{not} preserve eigenvalues, is upgraded to a similarity, and it is the reason the theorem must insist on orthonormal bases.
\end{proof}

% Extractor: Opus 5
\begin{exercise}[The Transition Matrix Between Two Orthonormal Bases]
\label{exc:transition_between_orthonormal_bases}
This is the step marked \qt{exc.} in the proof above, and it is what makes the whole argument work: it converts the congruence $Q\transp [B]_{\mathcal{D}'} Q$ of \cref{lem:change_of_basis_bilinear} into a similarity, so that eigenvalues are preserved. Let $V$ be a Euclidean space and let $\mathcal{D}, \mathcal{D}'$ be two orthonormal bases for $V$. Show that $[\id_V]_{\mathcal{D}'}^{\mathcal{D}} \in \Orth(n)$.
\exinfo{This exercise is taken from Prof.~Biran's handwritten lecture notes.}
\exsol{sol:transition_between_orthonormal_bases}
\end{exercise}
% Extractor: Gemini 3.1 Pro

% Thm. (27.a.7)
\begin{theorem}[Diagonalization of Symmetric Bilinear Forms]
\label{thm:diagonalization_symmetric_bilinear}
Let $V$ be a finite-dimensional vector space over a field $K$ of characteristic $0$. Let $B \colon V \times V \to K$ be a symmetric bilinear form. Then there exists a basis $\mathcal{C}$ of $V$ such that $[B]_{\mathcal{C}}$ is a diagonal matrix.
\end{theorem}

\begin{proof}
We proceed by induction on $n := \dim V$. We need to show there exists a basis $\mathcal{B} = (v_1, \dots, v_n)$ such that $B(v_i, v_j) = 0$ for all $i \neq j$.

If $B = 0$, or if $n = 1$, there is nothing to prove.

Let $n \ge 2$. Assume the statement of the theorem holds for all vector spaces of dimension $\le n-1$. Let $V$ be an $n$-dimensional vector space.
If $B(v, v) = 0$ for all $v \in V$, then the polarization formula (on \cpageref{thm:polarization_identity_bilinear}) shows $B = 0$ and we are done. (That identity is available here because $K$ has characteristic $0$, so in particular $2 \neq 0$ in $K$). So assume that there exists $v_1 \in V$ such that $B(v_1, v_1) \neq 0$.

Put $W := \Sp(v_1)$. The subspace $W$ is 1-dimensional (because $v_1 \neq 0$). Define the orthogonal complement with respect to $B$ as:
\[
W^{\perp_B} := \{ w \in V \mid B(v_1, w) = 0 \}.
\]
It is left as an exercise to show that $W^{\perp_B} \subseteq V$ is a subspace, see \cref{exc:orthogonal_complement_wrt_B}.

We claim that $W \cap W^{\perp_B} = \{0\}$. Indeed, if $w \in W \cap W^{\perp_B}$, then $w = c \cdot v_1$ and $B(v_1, w) = 0$. But $B(v_1, w) = B(v_1, c \cdot v_1) = c \cdot B(v_1, v_1)$. Since $B(v_1, v_1) \neq 0$, this implies $c = 0$, which means $w = 0$.

We claim that $V = W + W^{\perp_B}$.
Indeed, if $v \in V$, then the vector 
\[
u := v - \frac{B(v_1, v)}{B(v_1, v_1)} \cdot v_1
\]
belongs to $W^{\perp_B}$. To see this, we calculate:
\[
B(v_1, u) = B(v_1, v) - \frac{B(v_1, v)}{B(v_1, v_1)} \cdot B(v_1, v_1) = 0.
\]
This implies $u \in W^{\perp_B}$. So we can write $v$ as:
\[
v = \underbrace{\frac{B(v_1, v)}{B(v_1, v_1)} \cdot v_1}_{\in W} + \underbrace{u}_{\in W^{\perp_B}} \in W + W^{\perp_B}.
\]
It follows that $V = W \oplus W^{\perp_B}$. This implies $\dim W^{\perp_B} = n - 1$.

Now, $B|_{W^{\perp_B} \times W^{\perp_B}}$ is a symmetric bilinear form on an $(n-1)$-dimensional vector space ($W^{\perp_B}$). By the induction hypothesis, there exists a basis $\mathcal{C}' = (v_2, \dots, v_n)$ for $W^{\perp_B}$ such that $B(v_i, v_j) = 0$ for all $i \neq j$ with $i, j \ge 2$.

Take $\mathcal{C} := (v_1, \dots, v_n)$. Since $v_1 \in W$ and $v_2, \dots, v_n \in W^{\perp_B}$, we obtain that $B(v_i, v_j) = 0$ for all $i \neq j$.
\end{proof}

% Extractor: Opus 5
\begin{exercise}[The Orthogonal Complement with Respect to a Bilinear Form]
\label{exc:orthogonal_complement_wrt_B}
This is the step marked \qt{Exc.} in the proof above. With $B$, $V$ and $v_1$ as there, show that
\[ W^{\perp_B} := \{ w \in V \mid B(v_1, w) = 0 \} \]
is a linear subspace of $V$.
\exinfo{This exercise is taken from Prof.~Biran's handwritten lecture notes.}
\exsol{sol:orthogonal_complement_wrt_B}
\end{exercise}
% Extractor: Gemini 3.1 Pro

\begin{question}
What happens over $\mathbb{C}$?
\end{question}

% Thm. (27.a.8)
\begin{theorem}[Complex Symmetric Bilinear Forms]
\label{thm:complex_symmetric_bilinear}
Let $B \colon V \times V \to \mathbb{C}$ be a symmetric bilinear form on a vector space $V$ over $\mathbb{C}$, of dimension $n$. Then there exists a basis $\mathcal{C}$ of $V$ such that
\[
[B]_{\mathcal{C}} = \begin{pmatrix} I_r & 0 \\ 0 & 0 \end{pmatrix},
\]
for some $0 \le r \le n$. (The zero block in the bottom right has size $n - r$).

The number $r$ is independent of the basis $\mathcal{C}$: it depends only on $B$. In fact,
\[
n - r = \max \{ \dim W \mid W \subseteq V \text{ is a subspace such that } B|_{W \times W} \equiv 0 \}.
\]
The number $r$ is sometimes called the \newterm{rank} of $B$.
\end{theorem}

\begin{proof}
By the previous theorem (\cref{thm:diagonalization_symmetric_bilinear}), there exists a basis $\mathcal{C}' = (v'_1, \dots, v'_n)$ of $V$ such that
\[
[B]_{\mathcal{C}'} = \begin{pmatrix} \eta_1 & & 0 \\ & \ddots & \\ 0 & & \eta_n \end{pmatrix},
\]
where, after reordering the vectors of $\mathcal{C}'$ if necessary, $\eta_j \neq 0$ for all $1 \le j \le r$ and $\eta_j = 0$ for all $r < j \le n$. Here $r$ is simply the number of non-zero entries on the diagonal, so $0 \le r \le n$. (The reordering matters: the vectors $v'_j$ with $\eta_j \neq 0$ must come first, since we are about to divide by $\lambda_j$).

Choose $\lambda_j \in \mathbb{C}$ such that $\lambda_j^2 = \eta_j$ for $j = 1, \dots, r$. (This is always possible over $\mathbb{C}$, and $\lambda_j \neq 0$ because $\eta_j \neq 0$).
Define $v_j := \frac{1}{\lambda_j} v'_j$ for $j = 1, \dots, r$, and $v_q := v'_q$ for all $r < q \le n$.

Take $\mathcal{C} := (v_1, \dots, v_n)$. Then $B(v_i, v_j) = 0$ for all $i \neq j$. Furthermore, for $1 \le j \le r$, we have:
\[
B(v_j, v_j) = B\left(\frac{1}{\lambda_j} v'_j, \frac{1}{\lambda_j} v'_j\right) = \frac{1}{\lambda_j^2} \cdot \eta_j = 1,
\]
and for all $r+1 \le j \le n$, we have $B(v_j, v_j) = 0$.
\end{proof}

\subsection{More on Positive Definite Bilinear Forms}

Recall: A symmetric bilinear form $B \colon V \times V \to \mathbb{R}$ is called \newterm{positive definite} if $B(v, v) > 0$ for all $0 \neq v \in V$.

% Thm. (27.a.9)
\begin{theorem}[Sylvester's Criterion]
\label{thm:sylvesters_criterion}
Let $V$ be a finite-dimensional vector space over $\mathbb{R}$ and $B \colon V \times V \to \mathbb{R}$ a symmetric bilinear form. Then the following conditions are equivalent:
\begin{enumerate}[label=\textbf{(\alph*)}]
    \item $B$ is positive definite.
    \item There exists a basis $\mathcal{C} := (e_1, \dots, e_n)$ for $V$ such that the principal minors 
    \[ 
    \det(B(e_j, e_k))_{j,k=1, \dots, \ell} \quad 
    \]
    for all $\; 1 \le \ell \le n$, are positive.
    \item Like \textbf{(b)}, but instead of \qt{there exists a basis\dots}, it holds for \qt{all bases\dots}.
\end{enumerate}
\end{theorem}

\begin{proof}
\textbf{(b)} $\implies$ \textbf{(a)}: We proceed by induction on $n = \dim V$.

For $n = 1$, the assumption in \textbf{(b)} says $B(e_1, e_1) > 0$. This implies that for all $0 \neq c \in \mathbb{R}$, $B(c e_1, c e_1) = c^2 B(e_1, e_1) > 0$.

Suppose the implication holds for all spaces of dimension $\le n$. Let $V$ be an $(n+1)$-dimensional vector space, and $B \colon V \times V \to \mathbb{R}$ a symmetric bilinear form. Let $\mathcal{C} := (e_1, \dots, e_{n+1})$ be a basis for $V$ such that all the principal minors of $[B]_{\mathcal{C}}$ are positive.

Let $W := \Sp(e_1, \dots, e_n)$ with the basis $\mathcal{C}_W := (e_1, \dots, e_n)$.
The principal minors of $[B|_{W \times W}]_{\mathcal{C}_W}$ are exactly the same as the first $n$ principal minors of $[B]_{\mathcal{C}}$. Since these are positive, the induction hypothesis implies that $B|_{W \times W}$ is positive definite.

This implies that if we define $\langle w, w' \rangle_B := B(w, w')$, then $\langle \cdot, \cdot \rangle_B$ is a valid scalar product on $W$.

Consider the linear functional $L \colon W \to \mathbb{R}$ defined by $L(w) = B(w, e_{n+1})$.
By a previous result (Riesz Representation Theorem, \cref{thm:dual_space_inner_product_isomorphism}), there exists $w' \in W$ such that $L(w) = \langle w, w' \rangle_B$ for all $w \in W$. This implies:
\begin{equation}
\label{eq:sylvester_ortho}
B(w, w') = B(w, e_{n+1}) \quad \forall w \in W.
\end{equation}

Define a new basis for $V$, $\mathcal{C}' := (e_1, \dots, e_n, e_{n+1} - w')$. (It is left as an exercise to show this is a basis, see \cref{exc:shifted_family_is_basis}. Hint: $w' \in W$).
By \eqref{eq:sylvester_ortho}, we have $B(e_j, e_{n+1} - w') = 0$ for all $1 \le j \le n$. 

This implies that the matrix representation in the new basis is block diagonal:
\[
[B]_{\mathcal{C}'} = \begin{pmatrix} C & 0 \\ 0 & d \end{pmatrix} = S\transp [B]_{\mathcal{C}} S,
\]
where $d = B(e_{n+1} - w', e_{n+1} - w')$, $S = [\id_V]_{\mathcal{C}}^{\mathcal{C}'}$, and $C = [B|_{W \times W}]_{\mathcal{C}_W}$.

Taking the determinant, we get:
\[
\det([B]_{\mathcal{C}'}) = \det(C) \cdot d = \det(S)^2 \cdot \det([B]_{\mathcal{C}}).
\]
We know $\det(C) > 0$ (it is a principal minor), $\det(S)^2 > 0$ (since $S$ is invertible), and $\det([B]_{\mathcal{C}}) > 0$ (by assumption).
This implies that $d > 0$.

Let $v \in V$. We can write $v = w + \alpha \cdot (e_{n+1} - w')$, with $w \in W$ and $\alpha \in \mathbb{R}$.
Recall that $\langle w, e_{n+1} - w' \rangle_B = 0$.
This implies:
\[
B(v, v) = B(w, w) + \alpha^2 B(e_{n+1} - w', e_{n+1} - w') = B(w, w) + \alpha^2 d > 0,
\]
whenever either $w \neq 0$ or $\alpha \neq 0$. This implies that $B$ is positive definite.
The proof of \qt{\textbf{(b)} $\implies$ \textbf{(a)}} follows now by induction.

\textbf{(a)} $\implies$ \textbf{(b)}: Left as an exercise, see \cref{exc:sylvester_criterion_forward}. Similar ideas apply.
\end{proof}

% Extractor: Opus 5
\begin{exercise}[Shifting the Last Basis Vector Leaves a Basis]
\label{exc:shifted_family_is_basis}
This is the step Prof. Biran marks with \qt{exc.\ show this is a basis. Hint: $w' \in W$} in the proof above. Let $\mathcal{C} := (e_1, \dots, e_n, e_{n+1})$ be a basis for $V$ and let $w' \in W := \Sp(e_1, \dots, e_n)$. Show that $\mathcal{C}' := (e_1, \dots, e_n, e_{n+1} - w')$ is again a basis for $V$.
\exinfo{This exercise is taken from Prof.~Biran's handwritten lecture notes.}
\exsol{sol:shifted_family_is_basis}
\end{exercise}

\begin{exercise}[ of Sylvester's Criterion(\texorpdfstring{\cref{thm:sylvesters_criterion})]
\label{exc:sylvester_criterion_forward}
This is the implication Prof. Biran leaves with \qt{Exc. Similar ideas.} Prove that \textbf{(a)} $\implies$ \textbf{(b)} in \cref{thm:sylvesters_criterion}: if $B$ is positive definite, then the principal minors $\det(B(e_j, e_k))_{j,k = 1, \dots, \ell}$ are positive for $1 \le \ell \le n$.
\exinfo{This exercise is taken from Prof.~Biran's handwritten lecture notes.}
\exsol{sol:sylvester_criterion_forward}
\end{exercise}

\begin{exercise}[Testing Positive Definiteness via Sylvester's Criterion]
  \label{exc:sylvester_criterion_matrix_tests}
  Determine which of the following real symmetric matrices are positive definite by computing their leading principal minors:
  \[
  A := \begin{pmatrix} 3 & 3 & 2 & 3 \\ 3 & 1 & 1 & 2 \\ 2 & 1 & 2 & 1 \\ 3 & 2 & 1 & 3 \end{pmatrix}, \qquad
  B := \begin{pmatrix} 6 & 3 & 4 \\ 3 & 7 & 3 \\ 4 & 3 & 8 \end{pmatrix}, \qquad
  C := \begin{pmatrix} 3 & 0 & -1 & 0 \\ 0 & 6 & 1 & 1 \\ -1 & 1 & 8 & 2 \\ 0 & 1 & 2 & 5 \end{pmatrix}.
  \]
  \exhint{Apply Sylvester's Criterion (\cref{thm:sylvesters_criterion}): a real symmetric matrix is positive definite if and only if all its leading principal minors $\Delta_k := \det(M_{1 \le i,j \le k})$ are strictly positive for $1 \le k \le n$.
  For $A$: $\Delta_1 = 3 > 0$, but $\Delta_2 = \det\begin{pmatrix} 3 & 3 \\ 3 & 1 \end{pmatrix} = 3 - 9 = -6 < 0$, so $A$ is not positive definite.
  For $B$: $\Delta_1 = 6 > 0$, $\Delta_2 = 42 - 9 = 33 > 0$, and $\det(B) = 6(56-9) - 3(24-12) + 4(9-28) = 282 - 36 - 76 = 170 > 0$, so $B$ is positive definite.
  For $C$: compute the successive leading principal minors $\Delta_1, \Delta_2, \Delta_3, \Delta_4 > 0$ to verify that $C$ is positive definite.}
  \exinfo{This exercise is Problem 3 of Problem Sheet 26, Linear Algebra II, Spring Semester 2026, ETH Zurich (Lecturer: Prof. P. Biran).}
\exsol{sol:sylvester_criterion_matrix_tests}
\end{exercise}

\begin{exercise}[Equivalent Characterizations of Positive Definiteness \faHeart]
  \label{exc:positive_definite_equivalent_characterizations}
  Let $A \in \M_{n \times n}(\mathbb{R})$ be a symmetric matrix. Prove that the following statements are equivalent:
  \begin{enumerate}[label=\textbf{(\Alph*)}]
    \item $A$ is positive definite, i.e., $v\transp A v > 0$ for all non-zero vectors $v \in \mathbb{R}^n \setminus \{0\}$;
    \item All eigenvalues of $A$ are strictly positive;
    \item There exists an invertible symmetric matrix $S \in \M_{n \times n}(\mathbb{R})$ such that $S^2 = A$ (a symmetric square root).
  \end{enumerate}
  \exhint{Prove (A) $\implies$ (B) $\implies$ (C) $\implies$ (A).
  (A) $\implies$ (B): If $(\lambda, v)$ is an eigenvalue-eigenvector pair with $v \neq 0$, then $0 < v\transp A v = v\transp (\lambda v) = \lambda \|v\|^2 \implies \lambda > 0$.
  (B) $\implies$ (C): By the Spectral Theorem for symmetric matrices (\cref{thm:spectral_thm_matrices}), $A = O D O\transp$ with $O \in \Orth(n)$ and $D = \operatorname{diag}(\lambda_1, \dots, \lambda_n)$ where each $\lambda_i > 0$. Set $\sqrt{D} := \operatorname{diag}(\sqrt{\lambda_1}, \dots, \sqrt{\lambda_n})$ and $S := O \sqrt{D} O\transp$. Then $S\transp = S$, $S^2 = O D O\transp = A$, and $\det(S) = \prod \sqrt{\lambda_i} > 0$, so $S$ is invertible.
  (C) $\implies$ (A): For any $v \ne 0$, $v\transp A v = v\transp S^2 v = v\transp S\transp S v = \|Sv\|^2$. Since $S$ is invertible and $v \ne 0$, $Sv \ne 0$, hence $\|Sv\|^2 > 0$.}
  \exinfo{This exercise is Problem 4 of Problem Sheet 26, Linear Algebra II, Spring Semester 2026, ETH Zurich (Lecturer: Prof. P. Biran).}
\exsol{sol:positive_definite_equivalent_characterizations}
\end{exercise}
% Extractor: Gemini 3.1 Pro

% Creator: Gemini 3.1 Pro
\subsection{Solutions to the Exercises}
\label{sec:solutions_bilinear_quadratic}

% Creator: Opus 5
\begin{exercisesolution}[Proof of \cref{exc:symmetric_matrix_symmetric_form}]
\label{sol:symmetric_matrix_symmetric_form}
Let $v, w \in K^n$. The product $w\transp B v$ is a $1 \times 1$ matrix, that is, a scalar, and therefore, it equals its own transpose. Using this together with $B\transp = B$, we compute
\[
    B(w, v) = w\transp B v = \left(w\transp B v\right)\transp = v\transp B\transp w = v\transp B w = B(v, w),
\]
which is precisely the symmetry of the bilinear form.
\end{exercisesolution}

% Creator: Gemini 3.1 Pro
\begin{exercisesolution}[Proof of \cref{exc:two_nonzero_four_nonzero}]
\label{sol:two_nonzero_four_nonzero}
Assume $2 \neq 0$ in a field $K$. This means $1 + 1 \neq 0$. If $4 = 0$, then $2 + 2 = 0$, which means $2 \cdot 2 = 0$. Since $K$ is a field, it has no zero divisors. Thus, $2 \cdot 2 = 0$ implies $2 = 0$, which contradicts our assumption. Therefore, $4 \neq 0$.
\end{exercisesolution}

\begin{exercisesolution}[Proof of \cref{exc:orthogonal_complement_wrt_B}]
\label{sol:orthogonal_complement_wrt_B}
Let $w_1, w_2 \in W^{\perp_B}$ and $\alpha, \beta \in K$. We need to show that $\alpha w_1 + \beta w_2 \in W^{\perp_B}$. By definition of $W^{\perp_B}$, we have $B(v_1, w_1) = 0$ and $B(v_1, w_2) = 0$. Using the linearity of $B$ in the second argument, we compute:
\[
B(v_1, \alpha w_1 + \beta w_2) = \alpha B(v_1, w_1) + \beta B(v_1, w_2) = \alpha \cdot 0 + \beta \cdot 0 = 0.
\]
Thus, $\alpha w_1 + \beta w_2 \in W^{\perp_B}$, proving it is a subspace.
\end{exercisesolution}

\begin{exercisesolution}[Proof of \cref{exc:scaled_family_is_basis}]
\label{sol:scaled_family_is_basis}
We are given a basis $\mathcal{C}'' := (v''_1, \dots, v''_n)$. The new set $\mathcal{C} := (v_1, \dots, v_n)$ is formed by scaling each vector in $\mathcal{C}''$ by a non-zero scalar $c_i$ (where $c_i = 1/\lambda_i$ or $c_i = 1$). Since scaling basis vectors by non-zero scalars preserves linear independence and the span, $\mathcal{C}$ remains a basis for $V$.
\end{exercisesolution}

% Creator: Opus 5
\begin{exercisesolution}[Proof of \cref{exc:sylvester_needs_symmetry}]
\label{sol:sylvester_needs_symmetry}
Suppose that $\mathcal{C}$ is a basis for $V$ in which $M := [B]_{\mathcal{C}}$ has the form displayed in \cref{thm:sylvesters_inertia}. That matrix is symmetric, $M\transp = M$. Now let $v, w \in V$ be arbitrary and let $c, d \in \mathbb{R}^n$ be their coordinate vectors with respect to $\mathcal{C}$. By \cref{lem:matrix_rep_bilinear_form}, $B(v, w) = c\transp M d$, and since a $1 \times 1$ matrix equals its own transpose,
\[
    B(w, v) = d\transp M c = \left(d\transp M c\right)\transp = c\transp M\transp d = c\transp M d = B(v, w).
\]
So $B$ is symmetric. In other words, symmetry is not an artefact of the proof but a genuine necessary condition: no non-symmetric form can be brought to that shape in any basis.

For a form where the conclusion fails, take $V := \mathbb{R}^2$ and
\[ B(v, w) := v_1 w_2 - v_2 w_1, \]
which is bilinear with $B(e_1, e_2) = 1$ and $B(e_2, e_1) = -1$, hence not symmetric. By the argument just given, $[B]_{\mathcal{C}}$ never has the form of \cref{thm:sylvesters_inertia}. One can also see the failure directly: $q_B(v) = B(v, v) = v_1 v_2 - v_2 v_1 = 0$ for every $v$, so the associated quadratic form vanishes identically even though $B \neq 0$, and the numbers $k$ and $\ell$ could not possibly describe $B$.
\end{exercisesolution}

\begin{exercisesolution}[Proof of \cref{exc:transition_between_orthonormal_bases}]
\label{sol:transition_between_orthonormal_bases}
Write $\mathcal{D} = (d_1, \dots, d_n)$ and $\mathcal{D}' = (d'_1, \dots, d'_n)$, and put $Q := [\id_V]_{\mathcal{D}'}^{\mathcal{D}}$. By the definition of a transition matrix, column $j$ of $Q$ holds the coordinates of $d_j$ in the basis $\mathcal{D}'$, that is,
\[ d_j = \sum_{p=1}^n Q_{p,j} \, d'_p \qquad \text{for all } 1 \le j \le n. \]
Since $\mathcal{D}$ is orthonormal, $\langle d_j, d_k \rangle = \delta_{jk}$, and since $\mathcal{D}'$ is orthonormal, $\langle d'_p, d'_q \rangle = \delta_{pq}$. Expanding by bilinearity,
\[
    \delta_{jk} = \langle d_j, d_k \rangle = \left\langle \sum_{p=1}^n Q_{p,j} d'_p, \; \sum_{q=1}^n Q_{q,k} d'_q \right\rangle = \sum_{p,q=1}^n Q_{p,j} Q_{q,k} \, \delta_{pq} = \sum_{p=1}^n Q_{p,j} Q_{p,k} = \left(Q\transp Q\right)_{j,k}.
\]
Hence $Q\transp Q = I_n$, which is exactly the statement that $Q \in \Orth(n)$. (Equivalently, the displayed identity says that the columns of $Q$ form an orthonormal system, which is \cref{def:orthogonal_unitary_matrices}).
\end{exercisesolution}

% Creator: Gemini 3.1 Pro
\begin{exercisesolution}[Proof of \cref{exc:shifted_family_is_basis}]
\label{sol:shifted_family_is_basis}
We are given the basis $\mathcal{C} = (e_1, \dots, e_n, e_{n+1})$. We define $\mathcal{C}' := (e_1, \dots, e_n, e_{n+1} - w')$, where $w' \in W = \Sp(e_1, \dots, e_n)$. We can write $w' = \sum_{i=1}^n \alpha_i e_i$. 
Suppose a linear combination of $\mathcal{C}'$ equals zero:
\[
\sum_{i=1}^n c_i e_i + c_{n+1}(e_{n+1} - w') = 0 \implies \sum_{i=1}^n (c_i - c_{n+1}\alpha_i) e_i + c_{n+1} e_{n+1} = 0.
\]
Because $\mathcal{C}$ is a basis, all coefficients must be zero. In particular, $c_{n+1} = 0$. Substituting this back into the other coefficients yields $c_i = 0$ for all $1 \le i \le n$. Thus, the vectors in $\mathcal{C}'$ are linearly independent and form a basis.
\end{exercisesolution}

\begin{exercisesolution}[Proof of \cref{exc:sylvester_criterion_forward}]
\label{sol:sylvester_criterion_forward}
Assume $B$ is positive definite. Let $\mathcal{C} = (e_1, \dots, e_n)$ be any basis for $V$. For any $1 \le \ell \le n$, consider the subspace $W_\ell := \Sp(e_1, \dots, e_\ell)$. Since $B$ is positive definite on $V$, its restriction $B|_{W_\ell \times W_\ell}$ is also positive definite. The matrix representation of this restriction in the basis $\mathcal{E}_\ell := (e_1, \dots, e_\ell)$ is exactly the $\ell \times \ell$ principal submatrix of $[B]_{\mathcal{C}}$, so it is the $\ell$-th principal minor that we must show is positive.

% Creator: Opus 5
Apply \cref{thm:sylvesters_inertia} to the symmetric bilinear form $B|_{W_\ell \times W_\ell}$ on $W_\ell$. It supplies a basis $\mathcal{F}$ of $W_\ell$ in which the matrix of the restriction is $I_k \oplus (-I_m) \oplus 0$ with $k + m \le \ell$. Here $k$ is the largest dimension of a subspace on which the form is positive definite; since the form is positive definite on all of $W_\ell$, which has dimension $\ell$, we get $k = \ell$, and therefore, $m = 0$ and there is no zero block. In other words,
\[ \left[B|_{W_\ell \times W_\ell}\right]_{\mathcal{F}} = I_\ell, \qquad \text{so} \qquad \det\left[B|_{W_\ell \times W_\ell}\right]_{\mathcal{F}} = 1 > 0. \]
Now pass back to the basis $\mathcal{E}_\ell$. By \cref{lem:change_of_basis_bilinear}, with $S := [\id_{W_\ell}]_{\mathcal{F}}^{\mathcal{E}_\ell}$,
\[ \left[B|_{W_\ell \times W_\ell}\right]_{\mathcal{E}_\ell} = S\transp \cdot \left[B|_{W_\ell \times W_\ell}\right]_{\mathcal{F}} \cdot S, \]
and taking determinants gives $\det\left[B|_{W_\ell \times W_\ell}\right]_{\mathcal{E}_\ell} = \det(S)^2 \cdot 1 > 0$, because $S$ is invertible. This is the $\ell$-th principal minor of $[B]_{\mathcal{C}}$, and $\ell$ was arbitrary, which proves \textbf{(b)}.

Note that the basis $\mathcal{C}$ was arbitrary throughout, so what we have actually proved is the stronger implication \textbf{(a)} $\implies$ \textbf{(c)}. Since \textbf{(c)} $\implies$ \textbf{(b)} is immediate, the cycle \textbf{(a)} $\implies$ \textbf{(c)} $\implies$ \textbf{(b)} $\implies$ \textbf{(a)} closes, and all three conditions of \cref{thm:sylvesters_criterion} are equivalent. The notes state the three conditions but only ever discuss the passage between \textbf{(a)} and \textbf{(b)}, leaving \textbf{(c)} unproved; this is what closes that gap.
\end{exercisesolution}

% Official Solution (Problem Sheet 26, Exercise 1)
\begin{exercisesolution}[Solution to \cref{exc:polarization_quadratic_forms_series26}]
\label{sol:polarization_quadratic_forms_series26}
Let $K$ be a field with $\operatorname{char}(K) \neq 2$, let $V$ be a vector space over $K$, and let $B \colon V \times V \to K$ be a symmetric bilinear form.
Using the bilinearity of $B$, we expand the quadratic form $q_B(v + w)$:
\begin{align*}
q_B(v + w) &= B(v + w, v + w) \\
&= B(v, v) + B(v, w) + B(w, v) + B(w, w) \\
&= q_B(v) + B(v, w) + B(w, v) + q_B(w).
\end{align*}
Since $B$ is symmetric, $B(w, v) = B(v, w)$, which simplifies the expression to:
\[
q_B(v + w) = q_B(v) + 2 B(v, w) + q_B(w).
\]
Rearranging terms:
\[
2 B(v, w) = q_B(v + w) - q_B(v) - q_B(w).
\]
Since $\operatorname{char}(K) \ne 2$, the scalar $2 \in K$ is invertible. Dividing by $2$ yields the polarization identity:
\[
B(v, w) = \frac{1}{2} \big( q_B(v + w) - q_B(v) - q_B(w) \big). \qedhere
\]
\exinfo{Official Solution (Problem Sheet 26, Exercise 1). Derives the polarization identity expressing symmetric bilinear forms in terms of quadratic forms in characteristic $\ne 2$.}
\end{exercisesolution}

% Official Solution (Problem Sheet 26, Exercise 3)
\begin{exercisesolution}[Solution to \cref{exc:sylvester_criterion_matrix_tests}]
\label{sol:sylvester_criterion_matrix_tests}
We apply Sylvester's Criterion (\cref{thm:sylvesters_criterion}), which states that a real symmetric matrix is positive definite if and only if all its leading principal minors $\Delta_k := \det(M_{1 \le i,j \le k})$ are strictly positive for $1 \le k \le n$.
\begin{enumerate}[label=\textbf{(\alph*)}]
    \item \textbf{Matrix $A$:}
    The first two leading principal minors are:
    \[
    \Delta_1 = 3 > 0, \qquad \Delta_2 = \det \begin{pmatrix} 3 & 3 \\ 3 & 1 \end{pmatrix} = 3(1) - 3(3) = 3 - 9 = -6 < 0.
    \]
    Since $\Delta_2 < 0$, $A$ is \textbf{not positive definite}.

    \item \textbf{Matrix $B$:}
    The leading principal minors are:
    \begin{align*}
    \Delta_1 &= 6 > 0, \\
    \Delta_2 &= \det \begin{pmatrix} 6 & 3 \\ 3 & 7 \end{pmatrix} = 42 - 9 = 33 > 0, \\
    \Delta_3 &= \det(B) = 6(56 - 9) - 3(24 - 12) + 4(9 - 28) = 6(47) - 3(12) + 4(-19) \\
    &= 282 - 36 - 76 = 170 > 0.
    \end{align*}
    Since all leading principal minors are strictly positive ($\Delta_1, \Delta_2, \Delta_3 > 0$), $B$ is \textbf{positive definite}.

    \item \textbf{Matrix $C$:}
    Computing the successive leading principal minors:
    \begin{align*}
    \Delta_1 &= 3 > 0, \\
    \Delta_2 &= \det \begin{pmatrix} 3 & 0 \\ 0 & 6 \end{pmatrix} = 18 > 0, \\
    \Delta_3 &= \det \begin{pmatrix} 3 & 0 & -1 \\ 0 & 6 & 1 \\ -1 & 1 & 8 \end{pmatrix} = 3(48 - 1) - 1(0 - (-6)) = 3(47) - 6 = 141 - 6 = 135 > 0, \\
    \Delta_4 &= \det(C) = 3 \cdot \det\begin{pmatrix} 6 & 1 & 1 \\ 1 & 8 & 2 \\ 1 & 2 & 5 \end{pmatrix} + (-1) \cdot \det\begin{pmatrix} 0 & 6 & 1 \\ -1 & 1 & 2 \\ 0 & 1 & 5 \end{pmatrix} \\
    &= 3\big(6(36) - 1(3) + 1(-6)\big) - \big(30 - 1\big) = 3(207) - 29 = 621 - 29 = 592 > 0.
    \end{align*}
    Since all four leading principal minors are strictly positive, $C$ is \textbf{positive definite}. \qedhere
\end{enumerate}
\exinfo{Official Solution (Problem Sheet 26, Exercise 3). Tests positive definiteness of three matrices of sizes 2, 3, and 4 via leading principal minors.}
\end{exercisesolution}

% Official Solution (Problem Sheet 26, Exercise 4)
\begin{exercisesolution}[Solution to \cref{exc:positive_definite_equivalent_characterizations}]
\label{sol:positive_definite_equivalent_characterizations}
Let $A \in \M_{n \times n}(\mathbb{R})$ be symmetric. We prove the cycle of implications \textbf{(A)} $\implies$ \textbf{(B)} $\implies$ \textbf{(C)} $\implies$ \textbf{(A)}.
\begin{itemize}
    \item \textbf{(A) $\implies$ (B):}
    Let $\lambda$ be an eigenvalue of $A$ with non-zero eigenvector $v \in \mathbb{R}^n \setminus \{0\}$.
    By positive definiteness of $A$:
    \[
    0 < v\transp A v = v\transp (\lambda v) = \lambda (v\transp v) = \lambda \|v\|^2.
    \]
    Since $v \ne 0$, $\|v\|^2 > 0$, and dividing by $\|v\|^2$ gives $\lambda > 0$.
    Thus all eigenvalues of $A$ are strictly positive.

    \item \textbf{(B) $\implies$ (C):}
    Since $A$ is real symmetric, the Spectral Theorem (\cref{thm:spectral_thm_matrices}) ensures there exists an orthogonal matrix $O \in \Orth(n)$ such that:
    \[
    A = O D O\transp, \qquad D = \operatorname{diag}(\lambda_1, \dots, \lambda_n).
    \]
    By hypothesis, each $\lambda_i > 0$. We define $\sqrt{D} := \operatorname{diag}(\sqrt{\lambda_1}, \dots, \sqrt{\lambda_n})$ and set:
    \[
    S := O \sqrt{D} O\transp.
    \]
    Then $S$ is symmetric ($S\transp = (O \sqrt{D} O\transp)\transp = O \sqrt{D}\transp O\transp = O \sqrt{D} O\transp = S$), and:
    \[
    S^2 = (O \sqrt{D} O\transp)(O \sqrt{D} O\transp) = O \sqrt{D} (O\transp O) \sqrt{D} O\transp = O (\sqrt{D})^2 O\transp = O D O\transp = A.
    \]
    Moreover, $\det(S) = \det(O)\det(\sqrt{D})\det(O\transp) = \prod_{i=1}^n \sqrt{\lambda_i} > 0$, so $S$ is invertible.

    \item \textbf{(C) $\implies$ (A):}
    Let $S \in \M_{n \times n}(\mathbb{R})$ be an invertible symmetric matrix such that $S^2 = A$.
    For any non-zero vector $v \in \mathbb{R}^n \setminus \{0\}$:
    \[
    v\transp A v = v\transp S^2 v = v\transp S\transp S v = (Sv)\transp (Sv) = \|Sv\|^2.
    \]
    Since $S$ is invertible and $v \ne 0$, we have $Sv \ne 0$.
    By positive definiteness of the Euclidean norm, $\|Sv\|^2 > 0$, so $v\transp A v > 0$, proving that $A$ is positive definite. \qedhere
\end{itemize}
\exinfo{Official Solution (Problem Sheet 26, Exercise 4). Proves cyclic equivalence between positive definiteness, positive eigenvalues, and positive symmetric square root factorization.}
\end{exercisesolution}

% Creator: Opus 5
\begin{ainote}
\textbf{Two inverted representation-matrix indices}, the error class this pass
exists to catch, both in transition matrices and both invisible to
`tools/check-repmatrix.pl' because neither stands inside a product.

In the proof of \cref{thm:sylvesters_inertia}, the orthogonal $P$ produced by the
Spectral Theorem was introduced as $[\id_V]_{\mathcal{C}''}^{\mathcal{C}'} = P$,
with the basis of the target on top. The parenthetical immediately after it says
$v''_j = \sum_i P_{ij} v'_i$, that is, column $j$ of $P$ holds the coordinates of
$v''_j$ in $\mathcal{C}'$, which is $[\id_V]_{\mathcal{C}'}^{\mathcal{C}''}$; and
$[B]_{\mathcal{C}''} = P\transp A P$ is what
\cref{lem:change_of_basis_bilinear} yields for that reading and no other. The
same swap occurs in the proof of \cref{thm:sylvesters_criterion}, where
$S = [\id_V]_{\mathcal{C}'}^{\mathcal{C}}$ was written for the matrix that the
identity $[B]_{\mathcal{C}'} = S\transp [B]_{\mathcal{C}} S$ forces to be
$[\id_V]_{\mathcal{C}}^{\mathcal{C}'}$. Both are corrected;
\cref{lem:change_of_basis_bilinear} itself, the proof outline of
\cref{thm:sylvesters_euclidean} and the solution to
\cref{exc:sylvester_criterion_forward} were all checked and are right.

\textbf{Two results asserted without proof.} The Polarization Identity on
\cpageref{thm:polarization_identity_bilinear} carried none, although it is what
starts the induction of \cref{thm:diagonalization_symmetric_bilinear}; it is two
expansions and now has them. And \cref{lem:matrix_rep_bilinear_form} proves only
that a matrix representation \emph{exists}, then calls it \qt{the} matrix
representation; uniqueness is what entitles the later chapters to read
$([B]_{\mathcal{C}})_{jk} = B(v_j, v_k)$ off the name alone, as the proof of
\cref{lem:change_of_basis_bilinear} does in its first line, and it too is now
proved.

\textbf{One step whose point was the missing half.} The Remark before the
Conclusion shows that $\frac{1}{2}(B + B\transp)$ gives the same quadratic form
as $B$, without computing it and, more to the point, without saying that this
matrix is \emph{symmetric}, which is the whole content of the Conclusion that
follows. Both are supplied.

Steps that now name their reason: \cref{lem:char_orthogonal_matrices}
\textbf{(d)} for the three appearances of $P\transp = P^{-1}$;
the Proposition on \cpageref{prop:columns_of_representation_matrix} for what a
column of a transition matrix contains; \cref{lem:properties_char_poly} \textbf{(e)} with
\cref{aiprop:eigenvalue_iff_root} for \qt{similar, hence the same eigenvalues} in
\cref{thm:sylvesters_euclidean}, which is the exact point at which the congruence
of \cref{lem:change_of_basis_bilinear} has to be upgraded to a similarity, and so
the reason that theorem insists on orthonormal bases.

Flagged rather than changed: in the second item of the Example of standard forms
the letter $B$ denotes both the matrix and the bilinear form it defines, so that
the defining line reads $B(v, w) := v\transp B w$. Nothing is wrong, and the
notes do it too, but it is a trap for the eye of the same kind as the $P$ of ch.
17c; a parenthetical now says which $B$ is which.
\end{ainote}


